\chapter{Tiên đề của lí thuyết tập hợp}

\ 

Bạn đọc đã có thể từng học qua lí thuyết tập hợp từ cấp học phổ thông và được thực hiện xây dựng các tập hợp như
\begin{itemize}
   \item nhóm các nhân vật trong Đô-rê-mon: \{Nô-bi-ta; Xu-ka; Xê-kô; Chai-en; Đô-rê-mon; \dots\}, hay
   \item tập hợp các số tự nhiên: $\{0; 1; 2; 3; 4; \dots\}$.
\end{itemize}
Một định nghĩa tương đối ngây thơ có thể được đưa ra, cho phép một tập hợp có thể là một nhóm của các sự vật, hiện tượng cụ thể hay trừu tượng, yêu cầu duy nhất chúng ta cần thỏa mãn là có thể xác định được một cách rõ ràng một đối tượng có thuộc tập hợp đó hay không. Đây gọi là \defText{nguyên lý bao hàm không hạn chế}, viết dưới dạng ngôn ngữ lô-gích: Một tập hợp là một nhóm các hạng thức $x$ thỏa mãn $A(x)$ với $A$ là một vị từ chỉ phụ thuộc vào hạng thức tự do $x$. Không may, cách tiếp cận này dẫn đến những mâu thuẫn nội tại, nổi tiếng nhất là \defText{nghịch lý Rất-xen}\footnote{Bertrand Arthur William Russell (1872 -- 1970).}. Xét vị từ ``$x$ không thuộc vào tập hợp $x$''. Theo nguyên lí bao hàm, tồn tại một tập hợp $B$ bao gồm tất cả các tập hợp $x$ thỏa mãn $x \notin x$. Vậy, $B$ có thuộc $B$ hay không?
\begin{itemize}
   \item Nếu có, thì theo đúng điều kiện xác định của tập con, $B$ không thể tự chứa nó, suy ra $B \notin B$. Đây là một mâu thuẫn.
   \item Nếu không, thì vì $B$ không tự chứa chính nó, nó hoàn toàn thỏa mãn vị từ để trở thành một phần tử thuộc về tập hợp $B$, dẫn đến việc $B \in B$. Đây cũng là một mâu thuẫn!
\end{itemize}
Tình huống nan giải này khiến định nghĩa ngây thơ ban đầu bị phá sản, buộc các nhà toán học phải thiết lập một hệ tiên đề chặt chẽ và cẩn thận hơn nhằm ngăn chặn những mâu thuẫn như trên. Sau một quá trình phát triển, hệ tiên đề ZFC\footnote{Txéc-mê-lô -- Phren-ken (đặt tên theo Ernst Friedrich Ferdinand Zermelo (1871 -- 1953) và Abraham Fraenkel (1891 -- 1965)) cùng với tiên đề chọn (axiom of Choice).} đạt được nhiều thành công, và cũng sẽ là nền tảng mà chúng ta sẽ quan tâm.

\section{Đối tượng của lí thuyết tập hợp}

\ 

Từ cái tên, chúng ta cũng đã biết lí thuyết tập hợp sẽ tìm hiểu về tính chất của các \defText{tập hợp} (hay \defText{tập}). Chúng ta không định nghĩa tập hợp mà mặc nhiên cho nó tự tồn tại. Thêm vào đó, chúng ta cấp vị từ $\defMath{\in}\hspace{-0.25em}()$ nhận hạng thức là hai tập hợp $x$ và $y$ bất kì. Điều này có nghĩa $\defMath{\in}\hspace{-0.25em}\defMath{(x; y)}$ hay $\defMath{x \in y}$ là một mệnh đề có thể đúng hoặc sai. Nếu sai, chúng ta kí hiệu $\defMath{x \notin y}$. Nếu $x \in y$, chúng ta viết bằng lời ``$x$ thuộc $y$'' hay ``$x$ là một \defText{phần tử} của $y$''. Ngược lại, chúng ta viết ``$x$ không thuộc $y$''. 

Nếu như ghét người đọc, bạn đọc có thể viết theo thứ tự ngược lại:
\begin{align*}
    \defMath{y \ni x} &\iff \defMath{x \in y}; \\
    \defMath{y \not}\hspace{0.08em}\defMath{\ni x} &\iff \defMath{x \notin y}.
\end{align*}
Theo thứ tự, hai mệnh đề được phát biểu là ``$y$ chứa $x$'' và ``$y$ không chứa $x$''.

\section{Tiên đề ngoại diên}

\ 

Mặc dù cái tên có vẻ lạ, ý nghĩa của nó là một tính chất vô cùng quen thuộc của tập hợp: hai tập hợp cùng phần tử thì \defText{bằng nhau}. Chúng ta phát biểu nó như sau:
$$
\defMath{\forall x, \forall y, \left( x = y \iff  \forall z, \left( z \in x \iff z \in y \right) \right)}.
$$
Nếu không bằng nhau, kí hiệu $\defMath{x \neq y}$ được phát biểu là ``$x$ không bằng $y$''.
\begin{figure}[H]
   \begin{center}
      \begin{tikzpicture}[>=stealth]
         % Vẽ hai tập hợp
         \draw[thick] (-2,0) ellipse (1.2cm and 2cm);
         \node at (-2, 2.4) {Tập hợp $x$};
         
         \draw[thick] (2,0) ellipse (1.2cm and 2cm);
         \node at (2, 2.4) {Tập hợp $y$};
         
         % Các phần tử trong tập x
         \node[circle, fill, inner sep=1.2pt, label=left:$z_1$] (x1) at (-2, 1) {};
         \node[circle, fill, inner sep=1.2pt, label=left:$z_2$] (x2) at (-2, 0) {};
         \node[circle, fill, inner sep=1.2pt, label=left:$z_3$] (x3) at (-2, -1) {};
         
         % Các phần tử trong tập y
         \node[circle, fill, inner sep=1.2pt, label=right:$z_1$] (y1) at (2, 1) {};
         \node[circle, fill, inner sep=1.2pt, label=right:$z_3$] (y2) at (2, 0) {};
         \node[circle, fill, inner sep=1.2pt, label=right:$z_2$] (y3) at (2, -1) {};
         
         % Nối mũi tên hai chiều
         \draw[<->, thick, shorten >= 3pt, shorten <= 3pt] (x1) to[bend left=20] (y1);
         \draw[<->, thick, shorten >= 3pt, shorten <= 3pt] (x2) to (y3);
         \draw[<->, thick, shorten >= 3pt, shorten <= 3pt] (x3) to[bend left=10] (y2);
      \end{tikzpicture}
   \end{center}
   \caption{Minh họa hai tập hợp bằng nhau $x = y = \{z_1; z_2; z_3\}$}
\end{figure}
   

Trình bày bằng lời, chúng ta coi hai tập hợp là bằng nhau nếu mỗi phần tử trong một tập hợp thì đều thuộc tập hợp còn lại, và ngược lại. Trong trường hợp mà chúng ta chỉ có một chiều, thì chúng ta có định nghĩa \defText{tập con}:
$$
\defMath{\forall x, \forall y, \left( x \subseteq y \iff  \forall z, \left( z \in x \implies z \in y \right) \right)}.
$$

\begin{figure}[H]
   \begin{center}
      \begin{tikzpicture}[>=stealth]
         % Vẽ tập hợp x
         \draw[thick] (-2,0) ellipse (1.2cm and 2cm);
         \node at (-2, 2.4) {Tập hợp $x$};
         
         % Vẽ tập hợp y nghiêng và to hơn
         \draw[thick, rotate around={-20:(2.5,0)}] (2.5,0) ellipse (1.6cm and 2.5cm);
         \node at (2.5, 2.8) {Tập hợp $y$};
         
         % Các phần tử trong tập x
         \node[circle, fill, inner sep=1.2pt, label=left:$z_1$] (x1) at (-2, 1) {};
         \node[circle, fill, inner sep=1.2pt, label=left:$z_2$] (x2) at (-2, 0) {};
         \node[circle, fill, inner sep=1.2pt, label=left:$z_3$] (x3) at (-2, -1) {};
         
         % Các phần tử trong tập y (nhiều phần tử hơn, phân bố bên trong elip nghiêng)
         \node[circle, fill, inner sep=1.2pt, label=right:$z_1$] (y1) at (1.8, 1) {};
         \node[circle, fill, inner sep=1.2pt, label=right:$z_2$] (y2) at (1.8, -0.2) {};
         \node[circle, fill, inner sep=1.2pt, label=right:$z_3$] (y3) at (2.5, -1.2) {};
         \node[circle, fill, inner sep=1.2pt, label=right:$z_4$] (y4) at (3.2, 0.5) {};
         \node[circle, fill, inner sep=1.2pt, label=right:$z_5$] (y5) at (2.8, 1.6) {};
         
         % Nối mũi tên một chiều từ x sang y
         \draw[->, thick, shorten >= 3pt, shorten <= 3pt] (x1) to[bend left=15] (y1);
         \draw[->, thick, shorten >= 3pt, shorten <= 3pt] (x2) to (y2);
         \draw[->, thick, shorten >= 3pt, shorten <= 3pt] (x3) to[bend right=15] (y3);
      \end{tikzpicture}
   \end{center}
   \caption{Minh họa tập con $x \subseteq y$ hay $y \supseteq x$}
\end{figure}
Ngoài kí hiệu $\subseteq$, chúng ta còn có các định nghĩa sau: Với mọi tập hợp $x$ và $y$,
   \begin{itemize}
      \item $\defMath{x \supseteq y \iff y \subseteq x}$,
      \item $\defMath{x \subset y \iff x \subseteq y \land x \neq y}$\footnote{Trong nhiều tại liệu, $\subset$ có ý nghĩa tương đương với $\subseteq$.},
      \item $\defMath{x \subsetneq y \iff x \subset y}$,
      \item $\defMath{x \supset y \iff x \supseteq y \land x \neq y}$,
      \item $\defMath{x \supsetneq y \iff x \supset y}$.
   \end{itemize}

Nếu $x$ không phải là tập con của $y$ thì sao? Chúng ta có biến đổi như sau: Với mọi tập hợp $x$ và $y$,
\begin{align*}
   x \not \subseteq y &\iff \overline{\forall z, \left( z \in x \implies z \in y \right)}; \\
   x \not \subseteq y &\iff \exists z, \overline{\left( z \in x \implies z \in y \right)}; \\
   x \not \subseteq y &\iff \exists z, \left( z \in x \land z \not \in y \right).
\end{align*}
Tương tự với định nghĩa $\supseteq$, chúng ta cũng có $\forall x, \forall y, (x \not \supseteq y \iff y \not \subseteq x)$. Hình \ref{fig:li_thuyet_tap_hop:tien_de:khong_la_tap_con} minh họa cho $x \not \subseteq y$.

\begin{figure}[H]
   \begin{center}
      \begin{tikzpicture}[>=stealth]
         % Vẽ tập hợp x
         \draw[thick] (-2,0) ellipse (1.2cm and 2cm);
         \node at (-2, 2.4) {Tập hợp $x$};
         
         % Vẽ tập hợp y
         \draw[thick, rotate around={20:(2.5,0)}] (2.5,0) ellipse (1.6cm and 2.5cm);
         \node at (2.5, 2.8) {Tập hợp $y$};
         
         % Các phần tử trong tập x
         \node[circle, fill, inner sep=1.2pt, label=left:$z_1$] (x1) at (-2, 1) {};
         \node[circle, fill, inner sep=1.2pt, label=left:$z_2$] (x2) at (-2, 0) {};
         \node[circle, fill=red, inner sep=1.5pt, label={[red]left:$z_3$}] (x3) at (-2, -1) {};
         
         % Các phần tử trong tập y
         \node[circle, fill, inner sep=1.2pt, label=right:$z_1$] (y1) at (1.8, 1) {};
         \node[circle, fill, inner sep=1.2pt, label=right:$z_4$] (y2) at (1.8, 0) {};
         \node[circle, fill, inner sep=1.2pt, label=right:$z_2$] (y4) at (3.2, 0.5) {};
         \node[circle, fill, inner sep=1.2pt, label=right:$z_5$] (y5) at (2.8, -1.0) {};
         
         % Nối mũi tên một chiều từ x sang y cho z_1, z_2
         \draw[->, thick, shorten >= 3pt, shorten <= 3pt] (x1) to[bend left=15] (y1);
         \draw[->, thick, shorten >= 3pt, shorten <= 3pt] (x2) to[bend left=15] (y4);
         
         % Làm nổi bật rằng z không có tương ứng trong y
         \node[red] (text) at (0.25, -1) {$z_3 \notin y$};
         \draw[->, thick, red, densely dashed, shorten >= 2pt] (x3) -- (text);
      \end{tikzpicture}
   \end{center}
   \caption{Minh họa $x \not \subseteq y$ (tồn tại $z_3 \in x$ nhưng $z_3 \notin y$)}
   \label{fig:li_thuyet_tap_hop:tien_de:khong_la_tap_con}
\end{figure}

\section{Tiên đề tập rỗng}

\ 

Tiên đề này phát biểu rằng tồn tại một tập hợp mà không có phần tử nào. Trình bày dưới dạng kí hiệu:
$$
\defMath{\exists x, \forall y, y \not \in x}.
$$

Tập này được gọi là \defText{tập rỗng} và được kí hiệu là $\defMath{\emptyset}$. Tập này là duy nhất, do nếu giả sử tồn tại hai tập rỗng $\emptyset \neq \emptyset '$. Chúng ta có một số phân tích với phép $\neq$ như sau: Với $x$ và $y$ là hai tập hợp,
\begin{align*}
   x = y &\iff \forall z, (z \in x \iff z \in y); \\
   x \neq y &\iff \overline{\forall z, (z \in x \iff z \in y)}; \\
   x \neq y &\iff \exists z, \overline{(z \in x \iff z \in y)}.
\end{align*}
Sử dụng bảng giá trị chân lí, chúng ta có $\neg(P \iff Q) \iff (P \land \neg Q) \lor (\neg P \land Q)$.

\begin{table}[H]
	\centering
	\caption{Bảng giá trị chân lí của $\neg ( P \iff Q ) \iff ( P \land \neg Q ) \lor ( \neg P \land Q )$}
	\begin{tabular}{|c|c|cccccccccccccc|}
		\hline
		$P$ & $Q$ & $\neg$ & $(P$ & $\iff$ & $Q)$ & $\iff$ & $(P$ & $\land$ & $\neg$ & $Q)$ & $\lor$ & $(\neg$ & $P$ & $\land$ & $Q)$ \\
		\headerDivider
		Đ & Đ & S & Đ & Đ & Đ & \emphcolor{Đ} & Đ & S & S & Đ & S & S & Đ & S & Đ \\
		Đ & S & Đ & Đ & S & S & \emphcolor{Đ} & Đ & Đ & Đ & S & Đ & S & Đ & S & S \\
		S & Đ & Đ & S & S & Đ & \emphcolor{Đ} & S & S & S & Đ & Đ & Đ & S & Đ & Đ \\
		S & S & S & S & Đ & S & \emphcolor{Đ} & S & S & Đ & S & S & Đ & S & S & S \\
		\hline
	\end{tabular}
\end{table}
Sử dụng kết quả này để có
\begin{equation*}
   x \neq y \iff \exists z, (z \in x \land z \notin y) \lor (z \notin x \land z \in y).
\end{equation*}
Qua đó, từ $\emptyset \neq \emptyset '$, chúng ta cần phải có
$$
   \exists z, (z \in \emptyset \land z \notin \emptyset ') \lor (z \notin \emptyset \land z \in \emptyset ').
$$
Tuy nhiên, $z \in \emptyset$ luôn sai, và $z \in \emptyset '$ cũng luôn sai theo định nghĩa của tập rỗng. Do đó, cả hai phần của mệnh đề trên đều sai, dẫn đến một mâu thuẫn. Từ đó, chúng ta có chứng minh bằng phản chứng và mọi tập rỗng đều bằng nhau.

Nó có một tính chất mà có cảm giác như một nghịch lí: một tập rỗng là tập con của mọi tập hợp. Thật vậy, theo định nghĩa tập con,
$$
\forall y, (\emptyset \subseteq y \iff \forall z, (z \in \emptyset \implies z \in y)).
$$
Tuy nhiên, $z \notin \emptyset$ đúng với mọi $z$, cho nên vế giả thiết của $z \in \emptyset \implies z \in y$ luôn sai. Do đó, mệnh đề kéo theo này luôn đúng. Từ đó, hiển nhiên dẫn đến $\forall y, (\emptyset \subseteq y)$ đúng.

Liên quan một chút đến vấn đề triết học, người ta có thể đặt câu hỏi: ``Cần tối thiểu bao nhiêu phần tử cơ bản để tạo thành một thế giới?''. Với lí thuyết tập hợp, một tập rỗng là đủ. Trong những phần tiếp theo, bạn đọc sẽ thấy rằng tất cả các cấu trúc toán học đều có thể được xây dựng từ tập hợp. Tuy nhiên (và cũng tương đối hiển nhiên), một thế giới chỉ có một tập rỗng thì sẽ không có gì xảy ra cả, và cũng không có ai để nhận thức về nó. Do đó, mặc dù một tập rỗng là đủ để tạo ra một thế giới toán học, nhưng về mặt thực tế vật lí, chúng ta cần nhiều hơn thế.

Quay trở lại về các vấn đề liên quan đến tập hợp, chúng ta đã có điểm khởi đầu --- $\emptyset$. Những tiên đề tiếp theo sẽ cho phép xây dựng các tập hợp phức tạp hơn từ điểm khởi đầu này.

\section{Tiên đề cặp}

\ 

Tiên đề này cho phép chúng ta tạo ra một tập hợp mới từ hai tập hợp đã cho. Nó phát biểu rằng với mọi tập hợp $x$ và $y$, tồn tại một tập hợp $z$ sao cho $z$ chỉ chứa $x$ và $y$ làm phần tử của nó. Có thể thấy chỉ tồn tại một $z$ duy nhất theo tiên đề ngoại diên như chứng minh ở phần trước. Kí hiệu tập $z$ này là $\defMath{\{x; y\}}$. 

Dễ dàng thấy rằng $\{x; y\} = \{y; x\}$. Khi này, chúng ta có một \defText{đôi không thứ tự}. Một điểm nữa cần phải để ý, tiên đề không yêu cầu $x \neq y$, nên chúng ta hoàn toàn có thể xây dựng tập $\{x; x\}$ mà chỉ có một phần tử $x$. Đây là \defText{tập hợp đơn}, và có thể viết tắt tập này là $\defMath{\{x\}}$.